\chapter{ERP/CRM/SCM Integration in Commerce (Week 6)}

\section{Learning objectives}
\begin{itemize}
  \item Explain how commerce platforms integrate with ERP, CRM, and supply-chain systems across operational workflows.
  \item Identify systems of record for key entities such as customer, product, price list, order, invoice, and shipment.
  \item Analyze master data management, process ownership, and consistency risks in hybrid commerce architectures.
  \item Design integration flows for order-to-cash, returns, and B2B pricing using Odoo as a reference ERP environment.
\end{itemize}

\section{Why ERP/CRM/SCM integration is central to commerce}
Modern e-commerce experiences are often built in specialized commerce platforms, but commercial truth does not live only in the storefront.
Financial commitments, stock movements, supplier processes, customer account structures, invoicing, and back-office controls typically depend on enterprise systems such as ERP, CRM, and warehouse or supply-chain platforms.

This means a commerce platform cannot be architected only as a customer-facing product.
It must also participate in enterprise operating flows.
If that integration is weak, the result is familiar:
\begin{itemize}
  \item orders that appear confirmed online but fail in back-office processing,
  \item inventory figures that differ between storefront and warehouse,
  \item duplicate customer records across channels,
  \item B2B price lists that drift between sales and commerce systems,
  \item returns and refunds that create reconciliation problems in finance.
\end{itemize}

The core architectural task is not simply to ``connect Odoo'' or any ERP.
It is to define ownership, process boundaries, and trustworthy data movement across the enterprise stack.

\section{The enterprise systems perspective}
In a hybrid commerce architecture, enterprise systems typically serve different roles.

\subsection{ERP}
ERP systems coordinate structured operational and financial processes.
Using Odoo as a reference point, relevant modules often include:
\begin{itemize}
  \item \textbf{Sales:} quotations, sales orders, customer terms, sales workflows,
  \item \textbf{Inventory:} stock moves, reservations, warehouse operations, replenishment,
  \item \textbf{Accounting:} invoices, credit notes, payments, reconciliation,
  \item \textbf{CRM:} leads, opportunities, account history, commercial context.
\end{itemize}

\subsection{CRM}
CRM systems focus on relationship history, account management, sales interactions, lead qualification, and customer-service context.
In some organizations CRM and ERP overlap.
What matters architecturally is not the product label but the business responsibility assigned to the system.

\subsection{SCM and warehouse systems}
Supply-chain and warehouse platforms manage sourcing, supplier coordination, inbound logistics, warehouse execution, and delivery orchestration.
In smaller environments these functions may be embedded in ERP.
In larger environments they may be split into specialized systems such as WMS, TMS, or procurement platforms.

\section{Commerce versus enterprise ownership}
A recurring source of failure in commerce architecture is treating ownership as a technical detail rather than a business decision.
Teams must decide which domain owns each important entity and process state.

\subsection{Common ownership patterns}
\begin{itemize}
  \item The \textbf{commerce layer} often owns web sessions, carts, checkout state, personalization context, and channel-specific merchandising behavior.
  \item The \textbf{ERP layer} often owns invoices, accounting entries, stock movements, contractual customer records, and internal fulfillment commitments.
  \item The \textbf{CRM layer} may own lead status, sales activity, account segmentation, and service interactions.
  \item The \textbf{data platform} may own analytical aggregates and experimentation outputs but should not silently become the master for operational entities.
\end{itemize}

\subsection{Why ownership must be explicit}
If pricing is editable in both storefront and ERP, or if both commerce and CRM create customer records independently, inconsistencies are not accidental.
They are structural.
Integration design cannot fix a domain-ownership problem after the fact; it must start by making ownership explicit.

\section{Systems of record for key entities}
One of the most important artifacts in enterprise commerce architecture is a system-of-record matrix.
Representative entity choices include the following.

\subsection{Customer}
Customer identity is often split across systems:
\begin{itemize}
  \item the storefront may own channel account credentials and preferences,
  \item ERP may own bill-to and ship-to parties for financial processing,
  \item CRM may own relationship context, lead lifecycle, and account history.
\end{itemize}

The correct design is usually not ``one system owns everything.''
Instead, different aspects of customer identity are mastered in different places and linked through governed identifiers.

\subsection{Product}
Product ownership is also frequently split:
\begin{itemize}
  \item ERP may own procurement and internal stock-keeping information,
  \item a catalog or PIM layer may own rich digital merchandising content,
  \item analytics may derive behavioral product features.
\end{itemize}

The architecture must distinguish the commercial product definition from channel presentation and from analytical enrichment.

\subsection{Price list}
Price lists are especially sensitive in B2B environments.
Negotiated terms, tax logic, discount hierarchies, and approval rules often belong closer to ERP or enterprise sales processes than to storefront presentation logic.
If the storefront computes pricing independently while ERP maintains contractual prices, disputes are inevitable.

\subsection{Order}
The commerce platform often captures the initial order intent, but the enterprise question is whether ERP becomes the authoritative record for downstream operational and financial lifecycle.
This must be decided explicitly.
Some organizations treat commerce as the order system of record and export to ERP.
Others treat the storefront order as a pre-order until ERP accepts and materializes it.

\subsection{Invoice and payment reconciliation}
Invoices, ledger effects, and payment reconciliation usually belong in ERP or finance-owned systems.
Even if the storefront shows payment confirmation, the enterprise system remains authoritative for accounting truth.

\subsection{Shipment}
Shipment truth may belong in warehouse or logistics systems even if ERP or storefront mirrors the status.
The architecture should distinguish customer-facing delivery status from the internal system that owns shipment execution and exception handling.

\section{Master data management (MDM)}
Master data management is the discipline of defining and governing shared core entities across systems.
In commerce, the most common shared entities include customer, product, supplier, location, and price list.

\subsection{Why MDM matters}
Without MDM, organizations accumulate:
\begin{itemize}
  \item duplicate customer records created by multiple channels,
  \item product records with inconsistent category or attribute definitions,
  \item seller or supplier entities that differ across procurement and commerce tools,
  \item incompatible codes used by analytics, ERP, and storefront systems.
\end{itemize}

\subsection{Practical MDM principles}
\begin{itemize}
  \item Define authoritative identifiers and mapping rules.
  \item Separate creation rights from consumption rights.
  \item Document which attributes are mastered where.
  \item Maintain reconciliation workflows for conflicts and duplicates.
  \item Treat data stewardship as an operating responsibility, not an afterthought.
\end{itemize}

MDM is often perceived as bureaucratic.
In reality, it is a scaling mechanism for preventing repeated operational confusion.

\section{Order-to-cash in integrated commerce}
Order-to-cash is one of the most important enterprise commerce value streams because it crosses customer experience, logistics, and finance.

\subsection{Typical flow}
A simplified integrated order-to-cash process may include:
\begin{enumerate}
  \item customer places an order in the storefront,
  \item payment is authorized or payment terms are validated,
  \item order is persisted in commerce services,
  \item order is exported or synchronized to ERP,
  \item inventory is reserved or confirmed,
  \item picking, packing, and shipment are triggered,
  \item invoice is issued,
  \item payment is captured or reconciled,
  \item order and fulfillment statuses propagate back to customer-facing systems.
\end{enumerate}

\subsection{Architectural decisions inside the flow}
Key design choices include:
\begin{itemize}
  \item whether the storefront can confirm an order before ERP acceptance,
  \item whether stock is reserved in real time or reconciled later,
  \item whether ERP or commerce owns status transitions,
  \item how payment confirmation maps to invoicing and revenue recognition,
  \item which events trigger downstream systems such as warehouse or support tools.
\end{itemize}

\subsection{Failure modes}
Common failure modes include:
\begin{itemize}
  \item the order is captured online but fails to post in ERP,
  \item payment succeeds but the order export fails,
  \item the storefront shows available stock that has already been consumed elsewhere,
  \item fulfillment completes but shipment status never propagates back to customer channels,
  \item invoice creation and refund processing become inconsistent during cancellations.
\end{itemize}

These are not edge cases.
They are the expected complexity of distributed commercial workflows.

\section{Inventory, fulfillment, and availability}
Inventory integration is especially difficult because availability is both operationally critical and highly time-sensitive.

\subsection{Inventory concepts that should not be conflated}
\begin{itemize}
  \item \textbf{On-hand inventory:} physical stock currently available in storage.
  \item \textbf{Reserved inventory:} stock allocated to pending commitments.
  \item \textbf{Available-to-promise:} stock that can still be sold considering current commitments and business rules.
  \item \textbf{Future availability:} stock expected from replenishment or transfer.
\end{itemize}

If the storefront treats all four as the same number, overselling and customer disappointment become likely.

\subsection{Synchronization strategies}
Common strategies include:
\begin{itemize}
  \item near-real-time events for stock adjustments,
  \item short-lived availability caches,
  \item reservation services for high-demand items,
  \item periodic reconciliation against warehouse truth.
\end{itemize}

The right design depends on volume, latency needs, and oversell tolerance.
Enterprise architecture should make these tradeoffs explicit rather than hiding them behind ``inventory sync'' language.

\section{Returns and refunds}
Post-purchase flows are often under-modeled in early architecture work, yet they are essential to both customer trust and accounting correctness.

\subsection{Typical return flow}
A return process may involve:
\begin{enumerate}
  \item customer initiates a return,
  \item policy eligibility is checked,
  \item return merchandise authorization is created,
  \item item is received or inspected,
  \item stock is adjusted,
  \item refund or credit note is issued,
  \item customer and finance systems are updated.
\end{enumerate}

\subsection{Architectural risks}
Returns create risks when:
\begin{itemize}
  \item the storefront approves a return that finance policy rejects,
  \item inventory is restocked before inspection rules are applied,
  \item refunds are triggered twice through retries or manual intervention,
  \item credit notes in ERP do not align with customer-facing refund status.
\end{itemize}

These flows require the same rigor as order creation, including idempotency, state control, and clear ownership.

\section{B2B pricing and account structures}
B2B commerce introduces additional integration complexity because customer relationships are not purely transactional.
They often involve negotiated pricing, approval chains, sales-assist workflows, credit limits, and account hierarchies.

\subsection{Price lists and negotiated terms}
B2B price lists may depend on:
\begin{itemize}
  \item customer account,
  \item contract terms,
  \item region,
  \item product family,
  \item volume thresholds,
  \item approval status.
\end{itemize}

These rules are frequently managed in ERP or enterprise sales processes.
The storefront may present the result, but it should not become a shadow source of truth for contractual pricing.

\subsection{Account hierarchies}
One B2B customer may represent:
\begin{itemize}
  \item a legal entity,
  \item several branches,
  \item multiple buyers,
  \item different shipping destinations,
  \item separate billing relationships.
\end{itemize}

Architecture must therefore support organization-level identity and policy, not only individual user accounts.

\section{Using Odoo as a reference ERP}
Odoo is useful as a reference platform because it illustrates how enterprise workflows span modules rather than living in a single monolithic screen.

\subsection{Illustrative Odoo-aligned mapping}
\begin{itemize}
  \item \textbf{Sales:} quotations, orders, commercial terms, B2B account workflows.
  \item \textbf{Inventory:} stock movements, reservations, warehouse updates.
  \item \textbf{Accounting:} invoices, payments, credit notes, reconciliation.
  \item \textbf{CRM:} opportunity context, customer engagement history, account management.
\end{itemize}

The point of using Odoo here is not vendor lock-in.
It is to give students a concrete enterprise reference for how commerce-facing systems interact with structured back-office processes.

\subsection{What should remain vendor-neutral}
Even when Odoo is the teaching reference, the following concepts remain general:
\begin{itemize}
  \item system-of-record analysis,
  \item event and API contracts,
  \item master-data governance,
  \item process ownership,
  \item enterprise reconciliation and auditability.
\end{itemize}

\section{Integration risks and mitigation strategies}
Representative risks in ERP/CRM/SCM integration include:
\begin{itemize}
  \item \textbf{Duplicate sources of truth:} mitigated by explicit ownership and reconciliation.
  \item \textbf{Latency mismatch:} mitigated by choosing appropriate sync versus async boundaries.
  \item \textbf{Identifier mismatch:} mitigated by governed master identifiers and mapping tables.
  \item \textbf{Status divergence:} mitigated by explicit lifecycle models and state-mapping rules.
  \item \textbf{Operational blind spots:} mitigated by observability, error queues, and business-level monitoring.
\end{itemize}

Students should notice that these risks are architectural and organizational, not only technical.

\section{Artifacts for this chapter}
The expected architecture artifacts for ERP/CRM/SCM integration include:
\begin{itemize}
  \item an integration context diagram showing commerce, ERP, CRM, warehouse, payments, and analytics boundaries,
  \item a system-of-record matrix for shared entities,
  \item a small set of data contracts for critical events or APIs,
  \item and a short narrative for order-to-cash and returns ownership.
\end{itemize}

\section{Managerial implications}
Enterprise integration decisions affect more than system design.
They shape service reliability, customer trust, financial control, and the ability to scale channels.
If commerce and ERP are tightly coupled in every user interaction, customer experience becomes hostage to back-office latency.
If they are too loosely coupled without governance, the company loses commercial coherence.

The correct architectural aim is not maximal coupling or maximal decoupling.
It is controlled coordination with explicit ownership and recoverable failure handling.

\section{Multiple-choice questions (MCQs)}
\begin{enumerate}
  \item In an integrated commerce + ERP architecture, the \emph{system of record} is best defined as:
  \begin{enumerate}
    \item The system with the nicest UI
    \item The system that is most convenient for analytics
    \item The authoritative source of truth for a given entity (e.g., invoice) and its lifecycle
    \item The system that runs in the cloud
  \end{enumerate}

  \item Which flow is most strongly associated with \emph{order-to-cash}?
  \begin{enumerate}
    \item Creating a product image embedding
    \item Quote/order $\rightarrow$ delivery $\rightarrow$ invoice $\rightarrow$ payment reconciliation
    \item Building a recommendation model
    \item Defining a Kafka topic
  \end{enumerate}

  \item Master data management (MDM) is primarily concerned with:
  \begin{enumerate}
    \item Training deep learning models for search
    \item Defining and governing shared entities (customer, product, supplier) and preventing duplicates/inconsistencies
    \item Choosing the best database index
    \item Encrypting API requests
  \end{enumerate}

  \item A common integration risk when connecting commerce with ERP is:
  \begin{enumerate}
    \item Too many unit tests
    \item Duplicate sources of truth for prices, customers, or inventory
    \item Having an event schema
    \item Using idempotency keys
  \end{enumerate}

  \item In practice, why do integrations often require both APIs and events?
  \begin{enumerate}
    \item Because events are always cheaper than APIs
    \item Because some steps require immediate confirmation while others benefit from decoupling and retries
    \item Because ERP systems cannot expose APIs
    \item Because events guarantee zero duplicates without design effort
  \end{enumerate}
\end{enumerate}

\subsection*{Answer key}
\begin{enumerate}
  \item (c)
  \item (b)
  \item (b)
  \item (b)
  \item (b)
\end{enumerate}

\section{Exercises (short)}
\begin{enumerate}
  \item For each entity, choose a system of record (commerce vs Odoo) and justify: customer, product, price list, inventory on-hand, order, invoice, payment, return.
  \item Describe two failure modes of inventory sync and propose mitigations (e.g., eventual consistency, reservations, reconciliation jobs).
  \item Draft a ``data contract'' for an \texttt{OrderCreated} event: required fields, optional fields, and versioning rules.
\end{enumerate}

\section{Mini-case (Odoo-linked, vendor-neutral)}
\textbf{Scenario:} A hybrid company runs a headless storefront for B2C and a B2B portal for wholesale customers. Odoo is used for Sales, Inventory, Accounting, and CRM.

\textbf{Task:} Design the integration for three end-to-end processes:
\begin{itemize}
  \item \textbf{Order-to-cash:} from checkout to invoice posting and payment reconciliation.
  \item \textbf{Returns/refunds:} from return request to stock adjustment and credit note.
  \item \textbf{B2B pricing:} negotiated price lists and approvals without duplicating sources of truth.
\end{itemize}
For each process, state: key events, required synchronous calls, idempotency strategy, and the minimum monitoring you would implement.
